%% =====================================================================
%%  structures/informe-sr/structure.tex  --  INFORME DE SEMANA DE
%%  REPRESENTACION (ISR): cargador de la estructura editorial
%% ---------------------------------------------------------------------
%%  Concepto (pedido de Edison, 2026-09-03, v2): REVISTA INSTITUCIONAL
%%  DE ALTA GAMA. Menos elementos, mejor jerarquia: sin filetes, sin
%%  cajas ni lineas decorativas, sin franja textil, paleta minima
%%  (tinta + puka + gris; vino solo en la contraportada), prosa a dos
%%  columnas, pocas fotografias pero grandes, portada sobria.
%%
%%  Archivos (responsabilidad unica):
%%     structure.tex    paquetes, titulos, indice, cabecera/pie, sectores
%%     portada.tex      portada, apertura de articulo, membrete en linea,
%%                      separador (edicion interna), contraportada
%%     componentes.tex  kicker, cifra, cita, fotos, columnas, listas,
%%                      grafico; + fichas/cajas (solo edicion interna)
%%     tablas.tex       tablas (solo edicion interna: matriz, seguimiento)
%%     anexos.tex       anexos (solo edicion interna)
%%
%%  Se carga con \usarestructura{informe-sr} desde classes/informe-sr.cls.
%%  Preambulo del documento => nombres SIN @ y guardia anti-doble-carga.
%%  Iconos: PDF vectoriales en assets/iconos/ (Tabler): <nombre>-puka|tinta|blanco.pdf.
%% =====================================================================
\ifcsname despachoEstructuraInformeSR\endcsname
  \expandafter\endinput
\fi
\expandafter\def\csname despachoEstructuraInformeSR\endcsname{cargada}

%% --- Paquetes propios del tipo ----------------------------------------
\RequirePackage{tikz}
\usetikzlibrary{calc}
\RequirePackage{tcolorbox}
\tcbuselibrary{skins,breakable}
\RequirePackage{tabularray}
\UseTblrLibrary{booktabs}
\RequirePackage{pdfpages}
\RequirePackage{titlesec}
\RequirePackage{tocloft}
\RequirePackage{needspace}
\RequirePackage{ragged2e}
\RequirePackage{multicol}

%% --- Iconos (Tabler, PDF vectorial) -----------------------------------
\newcommand{\isricono}[3][puka]{\includegraphics[height=#3]{#2-#1.pdf}}
\newcommand{\isrdefsector}[3]{%
  \expandafter\def\csname isr-sec-ico-#1\endcsname{#2}%
  \expandafter\def\csname isr-sec-nom-#1\endcsname{#3}}
\isrdefsector{educacion}{school}{Educación}
\isrdefsector{salud}{building-hospital}{Salud}
\isrdefsector{infraestructura}{road}{Infraestructura}
\isrdefsector{agricultura}{plant}{Agricultura}
\isrdefsector{seguridad}{shield-check}{Seguridad}
\isrdefsector{justicia_ddhh}{scale}{Justicia y DD.\,HH.}
\isrdefsector{economia}{coin}{Economía}
\isrdefsector{gobierno}{building-community}{Gobierno y territorio}
\isrdefsector{comunicacion}{speakerphone}{Comunicación}
\isrdefsector{deporte}{ball-football}{Deporte}
\newcommand{\isrsectoricono}[3][puka]{\isricono[#1]{\csname isr-sec-ico-#2\endcsname}{#3}}
\newcommand{\isrsectornombre}[1]{\csname isr-sec-nom-#1\endcsname}

%% --- Tipografias de apoyo (una sola escala para todo el informe) -------
\newcommand{\isrtxkicker}{\fuentedisplaysemi\fontsize{7.6}{10}\selectfont\addfontfeatures{LetterSpace=13}}
\newcommand{\isrtxmeta}{\fuentedatos\fontsize{7.4}{9.5}\selectfont\color{colorisrtextosec}}
\newcommand{\isrtxpie}{\fuentedatos\fontsize{7}{9}\selectfont\color{colorisrtextosec}}
\newcommand{\isrtxcuerpo}{\fontsize{9.3}{11.8}\selectfont}
\newcommand{\isrtxcolumna}{\fontsize{9.9}{13.9}\selectfont}
\newcommand{\isrtxcolumnaancha}{\fontsize{9.9}{14.2}\selectfont}   % prosa a una columna: mismo cuerpo, interlinea apenas mayor

%% --- Titulo corrido y marca de capitulo -------------------------------
\newcommand{\isrtitulocorrido}{%
  \ifmetadef{titulocorrido}{\meta{titulocorrido}}{Informe de Semana de Representación}}
\renewcommand{\sectionmark}[1]{\markboth{\MakeUppercase{#1}}{}}
\renewcommand{\subsectionmark}[1]{}

%% --- Jerarquia de titulos (titlesec): sin numeros ni reglas ------------
%%  La jerarquia la dan el cuerpo, el peso y el aire, no las lineas.
%%  Los titulares NUNCA se parten con guion: a bandera y sin particion.
\newcommand{\isrsinguiones}{\hyphenpenalty=10000\exhyphenpenalty=10000\raggedright}
\setcounter{secnumdepth}{0}
\titleformat{\section}
  {\fuentedisplaynegra\fontsize{21}{24}\selectfont\color{colorisrtexto}\isrsinguiones}{}{0em}{}
\titlespacing*{\section}{0pt}{6pt}{14pt}
\titleformat{\subsection}
  {\fuentedisplaysemi\fontsize{11.5}{14}\selectfont\color{colorisrtexto}\isrsinguiones}{}{0em}{}
\titlespacing*{\subsection}{0pt}{14pt plus 3pt minus 2pt}{5pt}
\titleformat{\subsubsection}
  {\fuentedisplaysemi\fontsize{10}{13}\selectfont\color{colorisrprimario}}{}{0em}{}
\titlespacing*{\subsubsection}{0pt}{9pt plus 2pt}{3pt}
\newcommand{\seccion}[1]{\needspace{10\baselineskip}\section{#1}}
\newcommand{\subseccion}[1]{\needspace{8\baselineskip}\subsection{#1}}
\newcommand{\subsubseccion}[1]{\needspace{6\baselineskip}\subsubsection{#1}}

%% --- Indice (tocloft): limpio, sin puntos, folio a la derecha ----------
%%  OJO polyglossia fija «Índice general» al arrancar el documento y pisa
%%  cualquier \renewcommand del preambulo: el generador emite
%%  \renewcommand{\contentsname}{Contenido} en el CUERPO, justo antes de
%%  \tableofcontents (parte_creditos).
\setcounter{tocdepth}{1}
\renewcommand{\cfttoctitlefont}{\fuentedisplaynegra\fontsize{21}{24}\selectfont\color{colorisrtexto}}
\renewcommand{\cftaftertoctitle}{\par\vspace{2pt}}
\setlength{\cftbeforetoctitleskip}{0pt}
\setlength{\cftaftertoctitleskip}{12pt}
\renewcommand{\cftsecfont}{\fuentedisplaysemi\fontsize{10.5}{16}\selectfont\color{colorisrtexto}}
\renewcommand{\cftsecpagefont}{\fuentedatossemi\fontsize{9.5}{16}\selectfont\color{colorisrprimario}}
\renewcommand{\cftsecleader}{\hfill}
\setlength{\cftsecindent}{0pt}
\setlength{\cftsecnumwidth}{0pt}
\setlength{\cftbeforesecskip}{4pt}
\renewcommand{\cftsubsecfont}{\fontsize{9.5}{13}\selectfont\color{colorisrtextosec}}
\renewcommand{\cftsubsecpagefont}{\fuentedatos\fontsize{9}{13}\selectfont\color{colorisrtextosec}}
\renewcommand{\cftsubsecleader}{\hfill}
\setlength{\cftsubsecindent}{1.2em}
\setlength{\cftsubsecnumwidth}{0pt}
\setlength{\cftbeforesubsecskip}{2pt}

%% --- Columnas de prosa (revista) --------------------------------------
\setlength{\columnsep}{8mm}
\setlength{\columnseprule}{0pt}

%% --- Estilos de pagina: cabecera SOLO texto, sin filete ----------------
\newcommand{\isrcabeceraizq}{%
  {\fuentedatossemi\fontsize{6.4}{8}\selectfont\color{colorisrtextosec}%
    \addfontfeatures{LetterSpace=12}\MakeUppercase{\isrtitulocorrido}}}
\newcommand{\isrcabecerader}[1]{%
  {\fuentedisplaysemi\fontsize{6.4}{8}\selectfont\color{colorisrprimario}%
    \addfontfeatures{LetterSpace=10}#1}}
\newcommand{\isrpieizq}{%
  {\fuentedatos\fontsize{6.4}{8}\selectfont\color{colorisrtextosec}%
    \meta{autor}\ \textperiodcentered\ \meta{cargo}%
    \ifmetadef{distrito}{ por \meta{distrito}}{}}}
\newcommand{\isrpieder}{%
  {\fuentedatossemi\fontsize{8.5}{8.5}\selectfont\color{colorisrtexto}\thepage}}
\newcommand{\isranexomarca}{}

\fancypagestyle{isr}{%
  \fancyhf{}%
  \renewcommand{\headrulewidth}{0pt}%
  \renewcommand{\footrulewidth}{0pt}%
  \fancyhead[L]{\isrcabeceraizq}%
  \fancyhead[R]{\isrcabecerader{\leftmark}}%
  \fancyfoot[L]{\isrpieizq}%
  \fancyfoot[R]{\isrpieder}%
}
\fancypagestyle{isranexo}{%
  \fancyhf{}%
  \renewcommand{\headrulewidth}{0pt}%
  \renewcommand{\footrulewidth}{0pt}%
  \fancyhead[L]{\isrcabeceraizq}%
  \fancyhead[R]{\isrcabecerader{\isranexomarca}}%
  \fancyfoot[L]{\isrpieizq}%
  \fancyfoot[R]{\isrpieder}%
}
\pagestyle{isr}

%% --- Piezas de la estructura ------------------------------------------
%% =====================================================================
%%  structures/informe-sr/portada.tex  --  Portada, apertura de articulo,
%%  membrete en linea, separador (edicion interna) y contraportada
%% ---------------------------------------------------------------------
%%  Lo que va A SANGRE se dibuja como overlay TikZ anclado a la pagina
%%  fisica (remember picture => 2 pasadas, las de build.sh). Coordenadas
%%  en cm desde la esquina inferior izquierda del A4 (21.0 x 29.7).
%%
%%  Portada v2 (pedido de Edison 2026-09-03): SOBRIA. Fondo blanco, una
%%  sola fotografia a sangre horizontal (3:2, 21 x 14 cm), titulo grande,
%%  numero y periodo, ambito, identificacion del despacho y emblemas.
%%  Sin banda de color, sin franja textil, sin cifras, sin consigna.
%%
%%     \portadainforme[foto=, kicker=, titulo=, ambito=, codigo=]
%%     \articulo{kicker}{Titular}{Entrada}     apertura de cronica en flujo
%%     \membreteinforme                        membrete institucional en linea
%%     \capitulo[foto]{Titulo}{Subtitulo}      separador a pagina completa (edicion interna)
%%     \contraportadainforme
%% =====================================================================

%% --- Portada -----------------------------------------------------------
\pgfkeys{/isr/portada/.is family, /isr/portada,
  foto/.store in=\isrPortadaFoto,
  kicker/.store in=\isrPortadaKicker,
  titulo/.store in=\isrPortadaTitulo,
  ambito/.store in=\isrPortadaAmbito,
  codigo/.store in=\isrPortadaCodigo,
  default/.style={foto=, kicker=Informe de Semana de Representación,
    titulo={Semana de\\Representación}, ambito=, codigo=}}

\newcommand{\portadainforme}[1][]{%
  % OJO pgfkeys recorta UN solo espacio inicial por clave: sin espacio antes de #1.
  \pgfkeys{/isr/portada,default,#1}%
  \thispagestyle{empty}%
  \begin{tikzpicture}[remember picture, overlay]
    \begin{scope}[shift={(current page.south west)}]
      \fill[white] (0,0) rectangle (\paperwidth,\paperheight);
      % --- Linea institucional superior ------------------------------------
      \node[anchor=west, inner sep=0pt] at (2.0,28.30) {%
        \fuentedatossemi\fontsize{6.8}{8}\selectfont\color{colorisrtextosec}%
        \addfontfeatures{LetterSpace=16}CONGRESO DE LA REPÚBLICA · CÁMARA DE DIPUTADOS};
      \node[anchor=east, inner sep=0pt] at (19.0,28.30) {%
        \fuentedatossemi\fontsize{6.8}{8}\selectfont\color{colorisrtextosec}%
        \addfontfeatures{LetterSpace=16}PERIODO PARLAMENTARIO \meta{periodo}};
      % --- Kicker, titulo, ambito -----------------------------------------
      \node[anchor=north west, inner sep=0pt, text width=17cm, align=left] at (2.0,26.05) {%
        {\fuentedisplaysemi\fontsize{9.5}{12}\selectfont\color{colorisrprimario}%
          \addfontfeatures{LetterSpace=14}\MakeUppercase{\isrPortadaKicker}}\\[17pt]
        {\fuentedisplaynegra\fontsize{46}{49}\selectfont\color{colorisrtexto}\isrPortadaTitulo}\\[15pt]
        {\fontsize{13}{16}\selectfont\color{colorisrtextosec}\isrPortadaAmbito}};
      % --- Fotografia principal a sangre horizontal (21 x 14 cm) ------------
      \begin{scope}
        \clip (0,5.30) rectangle (\paperwidth,19.30);
        \if\relax\detokenize\expandafter{\isrPortadaFoto}\relax
          \fill[neutral200] (0,5.30) rectangle (\paperwidth,19.30);
          \node at (10.5,12.30) {%
            \fuentedatossemi\fontsize{7.5}{10}\selectfont\color{colorisrtextosec}%
            \addfontfeatures{LetterSpace=14}FOTOGRAFÍA PRINCIPAL
            \if\relax\detokenize\expandafter{\isrPortadaCodigo}\relax\else
              \ \textperiodcentered\ \isrPortadaCodigo\fi};
        \else
          \node[inner sep=0pt] at (10.5,12.30) {\includegraphics[width=\paperwidth]{\isrPortadaFoto}};
        \fi
      \end{scope}
      % --- Identificacion del despacho y emblemas ---------------------------
      \node[anchor=west, inner sep=0pt] at (2.0,3.10) {%
        \fuentedisplaynegra\fontsize{11.5}{14}\selectfont\color{colorisrtexto}%
        \addfontfeatures{LetterSpace=3}\MakeUppercase{\meta{autor}}};
      \node[anchor=west, inner sep=0pt] at (2.0,2.40) {%
        \fontsize{9.5}{12}\selectfont\color{colorisrtextosec}%
        \meta{cargo}\ifmetadef{distrito}{ por \meta{distrito}}{}};
      \node[anchor=east, inner sep=0pt] at (19.0,2.75) {%
        \includegraphics[height=1.35cm]{escudo-peru.png}\hspace{0.34cm}%
        \includegraphics[height=1.42cm]{logo-camara-mono.png}};
    \end{scope}
  \end{tikzpicture}%
  \null\clearpage
}

%% --- Apertura de cronica (articulo de revista, en flujo) ----------------
%%  \articulo{kicker}{Titular}{Entrada}: entra al indice y a los
%%  marcadores, fija la marca de cabecera. La foto de apertura, si la hay,
%%  se coloca ANTES con \isrfotoapertura (componentes.tex).
\newcommand{\articulo}[3]{%
  \par\Needspace{16\baselineskip}\vspace{14pt}%   % \Needspace (no \needspace): sin penalizacion blanda, que a una columna separaba foto y titular
  \phantomsection\addcontentsline{toc}{section}{#2}%
  \markboth{\MakeUppercase{#1}}{}%
  \noindent{\isrtxkicker\color{colorisrprimario}\MakeUppercase{#1}\par}%
  \vspace{7pt}%
  \noindent{\fuentedisplaynegra\fontsize{21}{24}\selectfont\color{colorisrtexto}\isrsinguiones#2\par}%
  \vspace{7pt}%
  \noindent{\fontsize{11.5}{15}\selectfont\itshape\color{colorisrtextosec}\isrsinguiones#3\par}%
  \vspace{10pt}}

%% --- Membrete institucional EN LINEA -----------------------------------
%%  Misma anatomia y colores que membretes/despacho (logos | caja negra
%%  del nombre | lema | filete negro), compuesto en el cuerpo para no
%%  ensanchar la cabecera de todo el documento. Solo en la presentacion.
\newcommand{\membreteinforme}{%
  \noindent\begin{minipage}{\linewidth}
    \begin{minipage}[c]{0.44\linewidth}\raggedright
      \includegraphics[height=1.50cm]{logo-camara-diputados.png}%
      \hspace{0.30cm}%
      \includegraphics[height=1.42cm]{escudo-peru.png}%
    \end{minipage}%
    \hfill
    \begin{minipage}[c]{0.54\linewidth}\raggedleft
      \setlength{\fboxsep}{0pt}%
      \colorbox{colormembrete}{%
        \begin{minipage}[c][1.30cm][c]{\despachoAnchoCaja}%
          \centering\color{white}%
          {\fuentemembretesemi\fontsize{10.5}{13}\selectfont\MakeUppercase{\meta{autor}}}\\[2.5pt]
          {\fuentemembrete\fontsize{9}{11}\selectfont\MakeUppercase{\meta{cargo}}}%
        \end{minipage}}\\[4.5pt]
      \makebox[\despachoAnchoCaja][c]{%
        {\itshape\fontsize{7.5}{9}\selectfont\color{neutral700}``\meta{lema}''}}%
    \end{minipage}\par\vspace{5pt}
    {\color{colorfilete}\rule{\linewidth}{0.6pt}}\par
  \end{minipage}\par\vspace{16pt}}

%% --- Separador de capitulo (pagina completa en vino) -- EDICION INTERNA --
\newcommand{\capitulo}[3][]{%
  \clearpage\thispagestyle{empty}%
  \phantomsection\addcontentsline{toc}{section}{#2}%
  \markboth{\MakeUppercase{#2}}{}%
  \begin{tikzpicture}[remember picture, overlay]
    \begin{scope}[shift={(current page.south west)}]
      \fill[colorisrprofundo] (0,0) rectangle (\paperwidth,\paperheight);
      \node[anchor=north west, inner sep=0pt, text width=15.5cm, align=left] at (2.0,22.0) {%
        {\fuentedisplaynegra\fontsize{28}{32}\selectfont\color{white}#2\par}
        \vspace{9pt}
        {\fontsize{12}{15.5}\selectfont\color{puka100}#3\par}};
      \if\relax\detokenize{#1}\relax\else
        \begin{scope}
          \clip (0,0) rectangle (\paperwidth,11.6);
          \node[inner sep=0pt] at (10.5,5.8) {\includegraphics[width=\paperwidth]{#1}};
        \end{scope}
      \fi
    \end{scope}
  \end{tikzpicture}%
  \null\clearpage
}

%% --- Contraportada -----------------------------------------------------
%%  Fondo vino (color profundo de la marca), retrato OFICIAL de cuerpo
%%  (\despachoRetratoCuerpo, lo fija el perfil del despacho), bloque de
%%  identificacion y contacto, emblemas. Sin franja textil ni consigna.
\newcommand{\contraportadainforme}{%
  \clearpage\thispagestyle{empty}%
  \begin{tikzpicture}[remember picture, overlay]
    \begin{scope}[shift={(current page.south west)}]
      \fill[colorisrprofundo] (0,0) rectangle (\paperwidth,\paperheight);
      \node[anchor=south west, inner sep=0pt] at (1.6,0) {%
        \includegraphics[height=13.2cm]{\despachoRetratoCuerpo}};
      %%  Nombre CORTO de marca (autorbanda) para que no se parta; el
      %%  nombre completo ya va en la portada y en el pie de cada pagina.
      \node[anchor=north west, inner sep=0pt, text width=9.0cm, align=left] at (9.95,17.2) {%
        {\fuentedatossemi\fontsize{7}{9}\selectfont\color{puka200}%
          \addfontfeatures{LetterSpace=16}CONGRESO DE LA REPÚBLICA\\CÁMARA DE DIPUTADOS\par}
        \vspace{10pt}
        {\fuentedisplaynegra\fontsize{22}{25}\selectfont\color{white}\isrsinguiones%
          \MakeUppercase{\ifmetadef{autorbanda}{\meta{autorbanda}}{\meta{autor}}}\par}
        \vspace{7pt}
        {\fuentedisplaysemi\fontsize{8.6}{12}\selectfont\color{puka200}%
          \addfontfeatures{LetterSpace=8}\MakeUppercase{\meta{cargo}\ifmetadef{distrito}{ POR \meta{distrito}}{}}\\
          PERIODO \meta{periodo}\par}
        \vspace{16pt}
        {\fontsize{9.5}{13}\selectfont\color{puka100}\meta{direccion}\par
          \vspace{3pt}\meta{correo}\ \textperiodcentered\ Cel.~\meta{celular}\par}
        \vspace{16pt}
        {\fontsize{8.8}{12}\selectfont\itshape\color{puka200}%
          \ifmetadef{notacontraportada}{\meta{notacontraportada}}{%
            Informe de la Semana de Representación elaborado por el
            despacho parlamentario y presentado conforme al Reglamento
            de la Cámara de Diputados.}\par}};
      \node[anchor=south east, inner sep=0pt] at (19.4,1.8) {%
        \includegraphics[height=1.40cm]{escudo-peru.png}\hspace{0.34cm}%
        \includegraphics[height=1.50cm]{logo-camara-blanco.png}};
    \end{scope}
  \end{tikzpicture}\null}

%% =====================================================================
%%  structures/informe-sr/componentes.tex  --  Componentes editoriales
%% ---------------------------------------------------------------------
%%  EDICION PUBLICA (revista): solo tipografia, aire y fotografia.
%%     \isrkicker{texto}                    etiqueta de seccion (sin regla)
%%     \isrcifra[ancho]{cifra}{etiqueta}{fuente}   dato grande (Inter)
%%     \isrcita{texto}{autor}{contexto}     cita destacada
%%     \isrfotoapertura{archivo}{pie}       foto grande de apertura (ancho de texto)
%%     \isrfotogrande{archivo}{pie}         foto a ancho de texto
%%     \isrfotosangre{archivo}{pie}         foto a sangre horizontal (21 cm)
%%     \begin{isrfotoytexto}{archivo}{pie} ... \end{isrfotoytexto}
%%                                          foto vertical + prosa al lado
%%     \isrfotocolumna{archivo}{pie}         foto al ancho de la columna, DENTRO
%%                                          de columnas (segunda foto de una cronica)
%%     \isrgaleriatitulo{texto}             kicker de la galeria del dia (pide sitio)
%%     \isrgaleriafila{a1}{pie1}{a2}{pie2}  fila indivisible de dos fotos 3:2 (a2 opcional)
%%     \isrgaleriatres{a1}{p1}{a2}{p2}{a3}{p3}  fila de tres (edicion a una columna)
%%     \isrfotoslot[ancho]{alto}{codigo}    espacio preparado, discreto
%%     \begin{columnas} ... \end{columnas}  prosa a dos columnas (una con \isrunacolumnatrue)
%%     \begin{isrpedidos} \item[P-011] ...  lista con codigo al margen
%%     \begin{isrcrono} \item[24 ago] ...   lista con fecha al margen
%%     \isrdato{etiqueta}{valor}            linea de ficha tecnica
%%     \isrgrafico[ancho]{archivo}{titulo}{fuente}
%%     \isrfuente{texto}
%%  EDICION INTERNA (documento de trabajo): fichas de actividad, cajas,
%%  chips (sin lineas verticales; pedido de Edison 2026-09-03).
%% =====================================================================

%% --- Kicker (sin regla) --------------------------------------------------
\newcommand{\isrkicker}[1]{%
  \par\needspace{4\baselineskip}\noindent
  {\isrtxkicker\color{colorisrprimario}\MakeUppercase{#1}}\par
  \vspace{5pt}}

%% --- Cifra (dato grande) ---------------------------------------------------
\newcommand{\isrcifra}[4][0.235\linewidth]{%
  \begin{minipage}[t]{#1}
    \setlength{\parskip}{0pt}\raggedright
    {\fuentedatossemi\fontsize{30}{32}\selectfont\color{colorisrprimario}#2\par}
    \vspace{4pt}
    {\fontsize{8.8}{11}\selectfont\color{colorisrtexto}#3\par}
    \if\relax\detokenize{#4}\relax\else
      \vspace{2pt}{\isrtxpie Fuente: #4\par}\fi
  \end{minipage}}

%% --- Linea de fuente ------------------------------------------------------
\newcommand{\isrfuente}[1]{\par\vspace{2pt}\noindent{\isrtxpie Fuente: #1\par}}

%% --- Cita destacada --------------------------------------------------------
\newcommand{\isrcita}[3]{%
  \par\needspace{6\baselineskip}\vspace{12pt}\noindent
  \begin{minipage}{\linewidth}
    \setlength{\parskip}{0pt}%
    {\fontsize{13.5}{17.5}\selectfont\itshape\color{colorisrprofundo}«#1»\par}
    \vspace{7pt}
    {\isrtxkicker\color{colorisrtextosec}\MakeUppercase{#2}%
      \if\relax\detokenize{#3}\relax\else\ \textperiodcentered\ #3\fi\par}
  \end{minipage}\par\vspace{12pt}}

%% --- Fotografias (sin marcos) ------------------------------------------------
\newcommand{\isrfotogrande}[2]{%
  \par\needspace{0.42\textheight}\vspace{8pt}\noindent
  \begin{minipage}{\linewidth}
    \setlength{\parskip}{0pt}%
    \includegraphics[width=\linewidth]{#1}\par\vspace{4pt}
    \if\relax\detokenize{#2}\relax\else{\isrtxpie #2\par}\fi
  \end{minipage}\par\vspace{12pt}}

%%  Apertura de cronica: exige casi la pagina completa libre; si no la
%%  hay, salta de pagina para que la foto abra arriba con su titular.
\newcommand{\isrfotoapertura}[2]{%
  \par\Needspace{0.85\textheight}\vspace{4pt}\noindent
  \begin{minipage}{\linewidth}
    \setlength{\parskip}{0pt}%
    \includegraphics[width=\linewidth]{#1}\par\vspace{4pt}
    \if\relax\detokenize{#2}\relax\else{\isrtxpie #2\par}\fi
  \end{minipage}\par}

%%  A sangre horizontal: la caja sobresale 2 cm por cada lado (margenes
%%  de 20 mm) y ocupa exactamente el ancho del papel. Va en el flujo.
\newcommand{\isrfotosangre}[2]{%
  \par\needspace{0.45\textheight}\vspace{8pt}\noindent
  \makebox[\linewidth][c]{\includegraphics[width=\paperwidth]{#1}}\par\vspace{4pt}
  \if\relax\detokenize{#2}\relax\else\noindent{\isrtxpie #2\par}\fi
  \vspace{12pt}}

%%  Foto vertical (2:3) a la izquierda + prosa a la derecha
\newenvironment{isrfotoytexto}[2]{%
  \par\needspace{0.5\textheight}\vspace{8pt}\noindent
  \begin{minipage}[t]{0.40\linewidth}\vspace{0pt}
    \setlength{\parskip}{0pt}%
    \includegraphics[width=\linewidth]{#1}\par\vspace{4pt}
    {\isrtxpie #2\par}
  \end{minipage}\hfill
  \begin{minipage}[t]{0.55\linewidth}\vspace{0pt}\isrtxcolumna\setlength{\parskip}{5pt}%
}{%
  \end{minipage}\par\vspace{12pt}}

%%  Foto de columna: va DENTRO de \begin{columnas}, al ancho de la columna,
%%  y fluye con la prosa (multicol la mueve a la columna siguiente si no
%%  cabe). Es la segunda foto de una cronica: la apertura sigue siendo la
%%  foto grande; con esta se admiten como maximo dos por cronica
%%  (pedido de Edison 2026-09-11). Sin marco, pie en \isrtxpie.
\newcommand{\isrfotocolumna}[2]{%
  \par\vspace{6pt}\noindent
  \begin{minipage}{\dimexpr\linewidth-\leftskip-\rightskip\relax}
    \setlength{\parskip}{0pt}%
    \includegraphics[width=\linewidth]{#1}\par\vspace{4pt}
    \if\relax\detokenize{#2}\relax\else{\isrtxpie #2\par}\fi
  \end{minipage}\par\vspace{10pt}}

%%  Galeria del dia (al menos una foto por actividad; pedido de Edison
%%  2026-09-11): filas de dos fotos al mismo formato 3:2, sin marco, con
%%  pie corto. Cada fila es una unidad indivisible que pide su espacio: si
%%  no cabe, pasa entera a la pagina siguiente. El segundo par puede ir
%%  vacio (numero impar de fotos). El titulo pide sitio para si y una fila,
%%  para no quedar huerfano al pie de pagina.
\newcommand{\isrgaleriatitulo}[1]{\par\Needspace{0.34\textheight}\vspace{10pt}\isrkicker{#1}}
\newcommand{\isrgaleriafila}[4]{%
  \par\Needspace{0.28\textheight}\vspace{4pt}\noindent
  \begin{minipage}[t]{0.49\linewidth}\vspace{0pt}
    \setlength{\parskip}{0pt}%
    \includegraphics[width=\linewidth]{#1}\par\vspace{3pt}
    \if\relax\detokenize{#2}\relax\else{\isrtxpie #2\par}\fi
  \end{minipage}\hfill
  \begin{minipage}[t]{0.49\linewidth}\vspace{0pt}
    \setlength{\parskip}{0pt}%
    \if\relax\detokenize{#3}\relax\else
      \includegraphics[width=\linewidth]{#3}\par\vspace{3pt}
      \if\relax\detokenize{#4}\relax\else{\isrtxpie #4\par}\fi
    \fi
  \end{minipage}\par\vspace{8pt}}

%%  Fila de tres: disponible, pero NO se usa en el ISR (Edison 2026-09-11: la
%%  galeria va siempre a dos columnas, nunca tres o mas; con dos fotos, grandes).
\newcommand{\isrgaleriatile}[3]{%
  \begin{minipage}[t]{#1}\vspace{0pt}
    \setlength{\parskip}{0pt}%
    \if\relax\detokenize{#2}\relax\else
      \includegraphics[width=\linewidth]{#2}\par\vspace{3pt}
      \if\relax\detokenize{#3}\relax\else{\isrtxpie #3\par}\fi
    \fi
  \end{minipage}}
\newcommand{\isrgaleriatres}[6]{%
  \par\Needspace{0.22\textheight}\vspace{4pt}\noindent
  \isrgaleriatile{0.315\linewidth}{#1}{#2}\hfill
  \isrgaleriatile{0.315\linewidth}{#3}{#4}\hfill
  \isrgaleriatile{0.315\linewidth}{#5}{#6}\par\vspace{8pt}}

%%  Espacio preparado (discreto): campo gris claro con etiqueta
\newcommand{\isrfotoslot}[3][\linewidth]{%
  \begin{minipage}[t]{#1}
    \begin{tcolorbox}[enhanced, colback=neutral200, boxrule=0pt, arc=0pt,
        height=#2, halign=center, valign=center, boxsep=0pt, left=4pt, right=4pt,
        before skip=0pt, after skip=0pt]
      {\fuentedatossemi\fontsize{6.4}{9}\selectfont\color{colorisrtextosec}%
        \addfontfeatures{LetterSpace=12}FOTOGRAFÍA PENDIENTE
        \if\relax\detokenize{#3}\relax\else\ \textperiodcentered\ #3\fi\par}
    \end{tcolorbox}
  \end{minipage}}

%% --- Columnas de prosa ----------------------------------------------------
%%  Dos columnas (revista) por defecto. Con \isrunacolumnatrue en el preambulo
%%  (edicion a una columna, pedido de Edison 2026-09-11) la prosa va a una sola
%%  columna con la medida reducida 1 cm por lado (≈ 90 caracteres) y la
%%  interlinea algo mayor; fotos de apertura y galerias siguen a ancho completo.
\newif\ifisrunacolumna
\newenvironment{columnasdos}{\begin{multicols}{2}\isrtxcolumna}{\end{multicols}}
\newenvironment{columnasuna}{\par\begingroup\isrtxcolumnaancha
  \setlength{\leftskip}{1cm}\setlength{\rightskip}{1cm}}{\par\endgroup}
\newenvironment{columnas}{\ifisrunacolumna\columnasuna\else\columnasdos\fi}%
  {\ifisrunacolumna\endcolumnasuna\else\endcolumnasdos\fi}

%% --- Listas con etiqueta al margen ------------------------------------------
\newlist{isrpedidos}{itemize}{1}
\setlist[isrpedidos]{label={}, leftmargin=1.45cm, labelwidth=1.15cm, labelsep=0.30cm,
  align=left, itemsep=4pt, topsep=4pt, parsep=0pt, labelindent=0pt,
  font=\fuentedatossemi\fontsize{7.6}{10}\selectfont\color{colorisrprimario}}
\newlist{isrcrono}{itemize}{1}
\setlist[isrcrono]{label={}, leftmargin=2.35cm, labelwidth=2.05cm, labelsep=0.30cm,
  align=left, itemsep=4pt, topsep=4pt, parsep=0pt, labelindent=0pt,
  font=\fuentedatossemi\fontsize{7.6}{10}\selectfont\color{colorisrprimario}}

%% --- Linea de ficha tecnica (etiqueta pequena + valor) ----------------------
\newcommand{\isrdato}[2]{%
  \par\noindent{\isrtxkicker\color{colorisrtextosec}\MakeUppercase{#1}\par}%
  \vspace{1.5pt}%
  {\fontsize{9.2}{12.5}\selectfont\color{colorisrtexto}#2\par}%
  \vspace{7pt}}

%% --- Grafico -------------------------------------------------------------------
\newcommand{\isrgrafico}[4][\linewidth]{%
  \par\needspace{10\baselineskip}\vspace{8pt}\noindent
  \begin{minipage}{#1}
    \setlength{\parskip}{0pt}%
    \isrkicker{#3}%
    \includegraphics[width=\linewidth]{#2}\par\vspace{3pt}
    {\isrtxpie Fuente: #4\par}
  \end{minipage}\par\vspace{12pt}}

%% =====================================================================
%%  EDICION INTERNA: chips, cajas y fichas de actividad
%% =====================================================================
\newcommand{\isrchip}[2][colorisrprimario]{%
  {\setlength{\fboxsep}{2.4pt}\colorbox{#1}{%
    \fuentedatossemi\fontsize{6.2}{6.2}\selectfont\color{white}%
    \addfontfeatures{LetterSpace=4}\,#2\,}}}
\newcommand{\isrchipclaro}[2][colorisrfila]{%
  {\setlength{\fboxsep}{2.4pt}\colorbox{#1}{%
    \fuentedatossemi\fontsize{6.2}{6.2}\selectfont\color{colorisrtexto}%
    \addfontfeatures{LetterSpace=4}\,#2\,}}}
\newcommand{\isrprioridad}[1]{%
  \ifcsname isr-prio-#1\endcsname\csname isr-prio-#1\endcsname\else\isrchipclaro{#1}\fi}
\expandafter\def\csname isr-prio-urgente\endcsname{\isrchip[puka700]{URGENTE}}
\expandafter\def\csname isr-prio-alta\endcsname{\isrchip[qillu600]{ALTA}}
\expandafter\def\csname isr-prio-ordinaria\endcsname{\isrchipclaro[neutral300]{ORDINARIA}}
\newcommand{\isrestado}[1]{%
  \ifcsname isr-est-#1\endcsname\csname isr-est-#1\endcsname\else\isrchipclaro{#1}\fi}
\expandafter\def\csname isr-est-registrado\endcsname{\isrchipclaro[neutral200]{REGISTRADO}}
\expandafter\def\csname isr-est-derivado\endcsname{\isrchipclaro[qillu100]{DERIVADO}}
\expandafter\def\csname isr-est-respondido\endcsname{\isrchipclaro[puka100]{RESPONDIDO}}
\expandafter\def\csname isr-est-cerrado\endcsname{\isrchip[exito]{CERRADO}}

\definecolor{qillu700}{HTML}{83600D}
\tcbset{
  isrcajabase/.style={enhanced, breakable, arc=0pt, outer arc=0pt,
    boxrule=0pt, left=10pt, right=10pt, top=8pt, bottom=8pt, boxsep=0pt,
    before skip=9pt, after skip=9pt,
    fontupper=\isrtxcuerpo, coltext=colorisrtexto,
    attach title to upper=\par\vspace{4pt},
    fonttitle=\isrtxkicker, coltitle=colorisrprimario},
  isrcaja-nota/.style={colback=colorisrfondo},
  isrcaja-alerta/.style={colback=colorisrfondosuave},
  isrcaja-compromiso/.style={colback=qillu50, coltitle=qillu700},
  isrcaja-dato/.style={colback=white, colframe=colorisrborde, boxrule=0.5pt},
  isrcaja-oscura/.style={colback=colorisrprofundo, coltext=white, coltitle=puka200},
}
\newtcolorbox{isrcaja}[2][nota]{isrcajabase, isrcaja-#1, title={\MakeUppercase{#2}}}

\newcommand{\isrfoto}[4][\linewidth]{%
  \begin{minipage}[t]{#1}
    \setlength{\parskip}{0pt}%
    \includegraphics[width=\linewidth]{#2}\par
    \vspace{3pt}%
    \if\relax\detokenize{#3}\relax\else
      {\fontsize{8.2}{10}\selectfont\color{colorisrtexto}#3\par}\fi
    \if\relax\detokenize{#4}\relax\else
      {\isrtxpie #4\par}\fi
  \end{minipage}}

\pgfkeys{/isr/ficha/.is family, /isr/ficha,
  id/.store in=\isrFichaId,
  titulo/.store in=\isrFichaTitulo,
  tipo/.store in=\isrFichaTipo,
  sector/.store in=\isrFichaSector,
  hora/.store in=\isrFichaHora,
  duracion/.store in=\isrFichaDuracion,
  lugar/.store in=\isrFichaLugar,
  ubicacion/.store in=\isrFichaUbicacion,
  participantes/.store in=\isrFichaParticipantes,
  resumenpedidos/.store in=\isrFichaResumenPedidos,
  documentos/.store in=\isrFichaDocumentos,
  foto/.store in=\isrFichaFoto,
  fotopie/.store in=\isrFichaFotoPie,
  fotoancho/.store in=\isrFichaFotoAncho,
  observacion/.store in=\isrFichaObs,
  default/.style={id=, titulo=, tipo=, sector=gobierno, hora=, duracion=,
    lugar=, ubicacion=, participantes=, resumenpedidos=, documentos=,
    foto=, fotopie=, fotoancho=0.34, observacion=}}
\tcbset{isrficha/.style={enhanced, breakable=false, colback=white,
  colframe=colorisrborde, boxrule=0.4pt, arc=0pt, outer arc=0pt,
  left=10pt, right=10pt, top=7pt, bottom=7pt, boxsep=0pt,
  before skip=9pt, after skip=9pt}}
\newcommand{\isrmeta}[2]{%
  \raisebox{-0.065cm}{\isricono[tinta]{#1}{0.29cm}}\hspace{2.5pt}#2\hspace{9pt}}
\newcommand{\isrfichacabecera}{%
  \noindent
  \isrchip{\isrFichaId}%
  \hspace{6pt}\raisebox{-0.1cm}{\isrsectoricono{\isrFichaSector}{0.36cm}}\hspace{3pt}%
  {\isrtxkicker\color{colorisrtextosec}\MakeUppercase{\isrsectornombre{\isrFichaSector}}}%
  \hfill
  {\isrtxkicker\color{colorisrtextosec}\MakeUppercase{\isrFichaTipo}}%
  \par\vspace{5pt}
  {\fuentedisplaysemi\fontsize{11.5}{13.5}\selectfont\color{colorisrtexto}\isrFichaTitulo\par}
  \vspace{3.5pt}
  {\isrtxmeta
    \if\relax\detokenize\expandafter{\isrFichaHora}\relax\else
      \isrmeta{clock}{\isrFichaHora
        \if\relax\detokenize\expandafter{\isrFichaDuracion}\relax\else\ (\isrFichaDuracion)\fi}\fi
    \if\relax\detokenize\expandafter{\isrFichaLugar}\relax\else
      \isrmeta{map-pin}{\isrFichaLugar
        \if\relax\detokenize\expandafter{\isrFichaUbicacion}\relax\else
          \ \textperiodcentered\ \isrFichaUbicacion\fi}\fi
    \if\relax\detokenize\expandafter{\isrFichaParticipantes}\relax\else
      \isrmeta{users}{\isrFichaParticipantes}\fi
    \par}
  \vspace{7pt}}
\newcommand{\isrfichapie}{%
  \par\vspace{6pt}
  \noindent{\isrtxmeta
    \if\relax\detokenize\expandafter{\isrFichaResumenPedidos}\relax\else
      \isrFichaResumenPedidos\fi
    \if\relax\detokenize\expandafter{\isrFichaDocumentos}\relax\else
      \hspace{10pt}\isrFichaDocumentos\fi
    \par}%
  \if\relax\detokenize\expandafter{\isrFichaObs}\relax\else
    \vspace{2.5pt}\noindent{\fontsize{7.6}{9.5}\selectfont\itshape\color{colorisrtextosec}%
      Nota de la fuente: \isrFichaObs\par}\fi}
\def\isrFotoSlot{slot}
\newenvironment{fichaactividad}[1]{%
  \pgfkeys{/isr/ficha,default,#1}%
  \begin{tcolorbox}[isrficha]
    \isrfichacabecera
    \if\relax\detokenize\expandafter{\isrFichaFoto}\relax
      \def\isrFichaTextoAncho{\linewidth}%
    \else
      \def\isrFichaTextoAncho{\dimexpr\linewidth-\isrFichaFotoAncho\linewidth-12pt\relax}%
    \fi
    \noindent\begin{minipage}[t]{\isrFichaTextoAncho}
      \setlength{\parskip}{3.5pt}\isrtxcuerpo\RaggedRight
}{%
    \end{minipage}%
    \if\relax\detokenize\expandafter{\isrFichaFoto}\relax\else
      \hfill
      \begin{minipage}[t]{\isrFichaFotoAncho\linewidth}
        \vspace{0pt}%
        \ifx\isrFichaFoto\isrFotoSlot
          \isrfotoslot{3.4cm}{\isrFichaId}%
        \else
          \isrfoto{\isrFichaFoto}{\isrFichaFotoPie}{}%
        \fi
      \end{minipage}%
    \fi
    \isrfichapie
  \end{tcolorbox}}
\newcommand{\isrcontexto}[1]{\noindent#1\par}
\newcommand{\isrsubtitulo}[1]{%
  \par\vspace{4pt}{\isrtxkicker\color{colorisrprimario}\MakeUppercase{#1}\par}\vspace{1.5pt}}
\newlist{isrhallazgos}{itemize}{1}
\setlist[isrhallazgos]{
  label={\textcolor{colorisrprimario}{\rule[0.09em]{0.40em}{0.40em}}},
  leftmargin=1.05em, labelsep=0.5em, labelindent=0pt,
  itemsep=1.5pt, topsep=2pt, parsep=0pt}

%% =====================================================================
%%  structures/informe-sr/tablas.tex  --  Tablas del informe (tabularray)
%% ---------------------------------------------------------------------
%%  Estilo de casa: cabecera en tinta con texto blanco (Archivo
%%  SemiBold, mayusculas), filas alternas en neutral200, filetes finos
%%  neutral300, cifras en Inter tabular alineadas a la derecha (regla
%%  dura de tokens/typography.json).
%%
%%  OJO tabularray NO expande macros ni al partir el cuerpo en celdas
%%  (filas con & y \\ literales: el generador las emite asi) ni dentro
%%  de la especificacion interior: por eso el estilo va ESCRITO en cada
%%  entorno y no en una macro compartida (fallo verificado 2026-09-03).
%%
%%     \begin{isrtabla}{colspec}   ... \end{isrtabla}       tabla corta
%%     \begin{isrtablalarga}{colspec} ... \end{isrtablalarga}  parte de pagina
%%     \begin{isrmatriz} fila & ... \\ \end{isrmatriz}       matriz de pedidos
%%     \begin{isrseguimiento} ... \end{isrseguimiento}       seguimiento
%%     \begin{isrcronograma} hora & \isrpunto & texto \\ \end{isrcronograma}
%%     \begin{isrfichatecnica} campo & valor \\ \end{isrfichatecnica}
%% =====================================================================

%% --- Textos de continuacion en espanol ---------------------------------
\DefTblrTemplate{contfoot-text}{default}{%
  \isrtxpie Continúa en la página siguiente}
\DefTblrTemplate{conthead-text}{default}{%
  \isrtxpie (continuación)}
\DefTblrTemplate{caption-tag}{default}{}
\DefTblrTemplate{caption-sep}{default}{}

%% --- Defaults comunes a TODAS las tablas del informe --------------------
\SetTblrInner[tblr,longtblr]{rowsep=3pt, colsep=4pt, stretch=0, vlines={0pt}}

%% --- Tabla corta con cabecera de casa -------------------------------------
\NewDocumentEnvironment{isrtabla}{m}
  {\par\vspace{4pt}\begin{tblr}{
      width=\linewidth, colspec={#1},
      row{1}={bg=colorisrtexto, fg=white, rowsep=4pt,
        font=\fuentedisplaysemi\fontsize{6.6}{8.5}\selectfont},
      row{2-Z}={font=\fontsize{8}{10}\selectfont},
      row{even}={bg=colorisrfila},
      hline{1}={0pt}, hline{2-Z}={0.3pt, colorisrborde}, hline{Z}={0.3pt, colorisrborde}}}
  {\end{tblr}\par\vspace{4pt}}

%% --- Tabla larga (se parte entre paginas, repite cabecera) ----------------
\NewDocumentEnvironment{isrtablalarga}{m}
  {\par\vspace{4pt}\begin{longtblr}[presep=0pt, postsep=0pt]{
      width=\linewidth, colspec={#1}, rowhead=1,
      row{1}={bg=colorisrtexto, fg=white, rowsep=4pt,
        font=\fuentedisplaysemi\fontsize{6.6}{8.5}\selectfont},
      row{2-Z}={font=\fontsize{8}{10}\selectfont},
      row{even}={bg=colorisrfila},
      hline{1}={0pt}, hline{2-Z}={0.3pt, colorisrborde}, hline{Z}={0.3pt, colorisrborde}}}
  {\end{longtblr}\par\vspace{4pt}}

%% --- Matriz de sistematizacion de pedidos --------------------------------
%%  Columnas: ID | Pedido o hallazgo | Entidad competente | Actuacion
%%  parlamentaria sugerida | Prioridad | Estado. El ID va en Inter puka.
\NewDocumentEnvironment{isrmatriz}{}
  {\par\vspace{4pt}\begin{longtblr}[presep=0pt, postsep=0pt]{
      width=\linewidth, rowhead=1,
      colspec={Q[l,wd=1.05cm] X[2.3,l] X[1.35,l] X[2.5,l] Q[c,wd=1.55cm] Q[c,wd=1.75cm]},
      row{1}={bg=colorisrtexto, fg=white, rowsep=4pt,
        font=\fuentedisplaysemi\fontsize{6.6}{8.5}\selectfont},
      row{2-Z}={font=\fontsize{7.6}{9.4}\selectfont},
      row{even}={bg=colorisrfila},
      cell{2-Z}{1}={font=\fuentedatossemi\fontsize{6.8}{9}\selectfont, fg=colorisrprimario},
      hline{1}={0pt}, hline{2-Z}={0.3pt, colorisrborde}, hline{Z}={0.3pt, colorisrborde}}
    ID & Pedido o hallazgo sistematizado & Entidad competente & Actuación parlamentaria sugerida & Prioridad & Estado \\}
  {\end{longtblr}\par\vspace{4pt}}

%% --- Tabla de seguimiento de compromisos ---------------------------------
\NewDocumentEnvironment{isrseguimiento}{}
  {\par\vspace{4pt}\begin{longtblr}[presep=0pt, postsep=0pt]{
      width=\linewidth, rowhead=1,
      colspec={Q[l,wd=1.35cm] X[2.6,l] X[1.3,l] Q[l,wd=1.9cm] Q[c,wd=1.55cm] Q[l,wd=2.1cm]},
      row{1}={bg=colorisrtexto, fg=white, rowsep=4pt,
        font=\fuentedisplaysemi\fontsize{6.6}{8.5}\selectfont},
      row{2-Z}={font=\fontsize{7.6}{9.4}\selectfont},
      row{even}={bg=colorisrfila},
      cell{2-Z}{1}={font=\fuentedatossemi\fontsize{6.8}{9}\selectfont, fg=colorisrprimario},
      hline{1}={0pt}, hline{2-Z}={0.3pt, colorisrborde}, hline{Z}={0.3pt, colorisrborde}}
    ID & Pedido & Entidad & Documento de salida & Estado & Responsable / plazo \\}
  {\end{longtblr}\par\vspace{4pt}}

%% --- Cronograma (linea de tiempo simple) ---------------------------------
\newcommand{\isrpunto}{\tikz[baseline=-0.55ex]{\fill[colorisrprimario] (0,0) circle (0.09cm);}}
\newcommand{\isrpuntoclaro}{\tikz[baseline=-0.55ex]{\draw[colorisrprimario, line width=0.6pt] (0,0) circle (0.09cm);}}
\NewDocumentEnvironment{isrcronograma}{}
  {\par\vspace{2pt}\begin{tblr}{
      width=\linewidth, colspec={Q[l,wd=1.75cm] Q[c,wd=0.35cm] X[l]},
      rowsep=2.5pt, colsep=3pt, hlines={0pt},
      column{1}={font=\fuentedatossemi\fontsize{7.8}{10}\selectfont, fg=colorisrprimario},
      column{3}={font=\fontsize{9}{11.5}\selectfont}}}
  {\end{tblr}\par\vspace{6pt}}
%%  Celda de actividad del cronograma: titulo + lugar en pequeno
\newcommand{\isrhitotexto}[2]{{\fuentedisplaysemi\fontsize{8.8}{11}\selectfont #1}\newline{\isrtxpie #2}}

%% --- Ficha tecnica del documento (pagina de creditos) --------------------
\NewDocumentEnvironment{isrfichatecnica}{}
  {\par\begin{tblr}{
      width=\linewidth, colspec={Q[l,wd=4.2cm] X[l]},
      rowsep=3.5pt, colsep=6pt,
      hlines={0.3pt, colorisrborde}, hline{1}={0pt},
      column{1}={font=\fuentedisplaysemi\fontsize{7.4}{10}\selectfont, fg=colorisrtextosec},
      column{2}={font=\fontsize{9}{11.5}\selectfont}}}
  {\end{tblr}\par}

%% =====================================================================
%%  structures/informe-sr/anexos.tex  --  Anexos documentales
%% ---------------------------------------------------------------------
%%  Cada anexo abre con una CARATULA propia (numero, titulo, descripcion,
%%  entrada en el indice y en los marcadores) y, si es un documento PDF
%%  (oficio, memorando, acta), se incluyen sus paginas escaladas con
%%  pdfpages conservando la cabecera corrida del informe (estilo isranexo,
%%  que nombra el anexo en curso).
%%
%%     \isranexoportada{N}{Titulo}{Descripcion}
%%     \isranexopdf[opciones pdfpages]{archivo.pdf}
%%     \isranexoimagen{archivo}{pie}          escaneo/imagen a pagina completa
%% =====================================================================

\newcommand{\isranexoportada}[3]{%
  \clearpage\thispagestyle{isr}%
  \gdef\isranexomarca{ANEXO #1 · \MakeUppercase{#2}}%
  \markboth{ANEXO #1}{}%
  \phantomsection\addcontentsline{toc}{subsection}{Anexo #1. #2}%
  \vspace*{2.6cm}
  \noindent{\isrtxkicker\color{colorisrprimario}ANEXO #1\par}
  \vspace{9pt}
  \noindent{\fuentedisplaynegra\fontsize{22}{26}\selectfont\color{colorisrtexto}#2\par}
  \vspace{7pt}
  \noindent{\color{colorisracento}\rule{2.4cm}{2pt}\par}
  \vspace{11pt}
  \noindent{\fontsize{10.5}{14}\selectfont\color{colorisrtextosec}#3\par}
}

\newcommand{\isranexopdf}[2][]{%
  \includepdf[pages=-, scale=0.88, offset=0 -3mm, frame=false,
    pagecommand={\thispagestyle{isranexo}}, #1]{#2}}

\newcommand{\isranexoimagen}[2]{%
  \clearpage\thispagestyle{isranexo}%
  \noindent\begin{minipage}{\linewidth}\centering
    \includegraphics[width=\linewidth, height=0.86\textheight, keepaspectratio]{#1}\par
    \vspace{4pt}{\isrtxpie #2\par}
  \end{minipage}\par}

